\subsection{求解与结果}
基准时间窗下，先在 1.4--4.8 的粗网格和 3.0--3.4 的加密网格上扫描 $\lambda_{\max}$，得到首个连续正区间出现在 $K=3.2186$ 附近，并在 3.22--3.40 区间保持为正。表 \ref{tab:q4robust} 给出的局部复核表明，在总积分时长固定为 90 的前提下，$K=3.20$ 与 $K=3.22$ 之间的符号跨越对相邻步长和微扰初值仍可保留。

\begin{table}[htbp]
\centering
\caption{问题四短时间窗局部复核}
\label{tab:q4robust}
\small
\begin{tabularx}{0.95\textwidth}{>{\raggedright\arraybackslash}p{4.1cm} >{\centering\arraybackslash}p{1.9cm} >{\centering\arraybackslash}p{1.9cm} >{\centering\arraybackslash}p{1.9cm} X}
\toprule
设置 & $\lambda_{\max}(3.20)$ & $\lambda_{\max}(3.22)$ & 判别 & 说明 \\
\midrule
基准步长 $dt=0.010$ & $-0.00135$ & $0.00117$ & 符号跨越 & 3.22 附近进入正指数区间 \\
较小步长 $dt=0.008$ & $-0.00134$ & $0.00117$ & 符号跨越 & 结果与基准设置一致 \\
较大步长 $dt=0.012$ & $-0.00129$ & $0.00123$ & 符号跨越 & 结果与基准设置一致 \\
微扰初值 $(0.8005,0.1995,0.2)$ & $-0.00087$ & $0.00165$ & 符号跨越 & 微扰后仍保留同样的符号关系 \\
\bottomrule
\end{tabularx}
\normalsize
\end{table}

积分时长拉长后，正指数区间明显左移。总时长为150时，3.0--3.4加密区间内各点均为正；总时长为200时，连续正区间从3.11开始，且$\lambda_{\max}(3.20)=0.00579$、$\lambda_{\max}(3.22)=0.00701$两点均为正。扩展到5组初值、5个步长和3个积分时长的75组组合后，仅有5组（6.7\%）保留$\lambda_{\max}(3.20)<0<\lambda_{\max}(3.22)$。这不能支持稳定临界点的解释，即该符号跨越在扩展设置下不具备统计稳健性。

基准短时间窗在$K=3.2186$附近出现由负转正的局部跨越；积分时长和数值设置改变后，跨越位置随之移动，因此$K=3.2186$不能作为唯一临界点。本文只保留“承载力提高会增强系统复杂性”这一方向性结论。若需精确定位，还须采用更长积分时长（$\ge 500$）并与分岔图交叉核对，图\ref{fig:q4}给出不同积分时长下的扫描结果。

\begin{figure}[htbp]
\centering
\includegraphics[width=0.92\textwidth]{fig04_q4_lyapunov_scan.png}
\caption{问题四不同积分时长下的有限时间李雅普诺夫指数扫描}
\label{fig:q4}
\end{figure}

